\documentclass[10pt, letterpaper]{article}

% Packages:
\usepackage[
    ignoreheadfoot, % set margins without considering header and footer
    top=0.9 cm, % seperation between body and page edge from the top
    bottom=0.7 cm, % seperation between body and page edge from the bottom
    left=1.8 cm, % seperation between body and page edge from the left
    right=1.8 cm, % seperation between body and page edge from the right
    footskip=1.0 cm, % seperation between body and footer
    % showframe % for debugging 
]{geometry} % for adjusting page geometry
\usepackage{titlesec} % for customizing section titles
\usepackage{tabularx} % for making tables with fixed width columns
\usepackage{array} % tabularx requires this
\usepackage[dvipsnames]{xcolor} % for coloring text
\definecolor{primaryColor}{RGB}{0, 79, 144} % define primary color
\usepackage{enumitem} % for customizing lists
\usepackage{fontawesome5} % for using icons
\usepackage{amsmath} % for math
\usepackage[
    pdftitle={Resume},
    pdfauthor={Student},
    pdfcreator={LaTeX with RenderCV},
    colorlinks=true,
    urlcolor=primaryColor
]{hyperref} % for links, metadata and bookmarks
\usepackage[pscoord]{eso-pic} % for floating text on the page
\usepackage{calc} % for calculating lengths
\usepackage{bookmark} % for bookmarks
\usepackage{lastpage} % for getting the total number of pages
\usepackage{changepage} % for one column entries (adjustwidth environment)
\usepackage{paracol} % for two and three column entries
\usepackage{ifthen} % for conditional statements
\usepackage{needspace} % for avoiding page brake right after the section title
\usepackage{iftex} % check if engine is pdflatex, xetex or luatex

% Ensure that generate pdf is machine readable/ATS parsable:
\ifPDFTeX
    \input{glyphtounicode}
    \pdfgentounicode=1
    % \usepackage[T1]{fontenc} % this breaks sb2nov
    \usepackage[utf8]{inputenc}
    \usepackage{lmodern}
\fi

% Some settings:
\AtBeginEnvironment{adjustwidth}{\partopsep0pt} % remove space before adjustwidth environment
\pagestyle{empty} % no header or footer
\setcounter{secnumdepth}{0} % no section numbering
\setlength{\parindent}{0pt} % no indentation
\setlength{\topskip}{0pt} % no top skip
\setlength{\columnsep}{0cm} % set column seperation
\makeatletter
\let\ps@customFooterStyle\ps@plain % Copy the plain style to customFooterStyle
\makeatother
\pagestyle{customFooterStyle}

\titleformat{\section}{\needspace{3\baselineskip}\bfseries\large}{}{0pt}{}[\vspace{1pt}\titlerule]

\titlespacing{\section}{
    % left space:
    1pt
}{
    % top space:
    0.05 cm
}{
    % bottom space:
    0.05 cm
} % section title spacing

\renewcommand\labelitemi{$\circ$} % custom bullet points
\newenvironment{highlights}{
    \begin{itemize}[
        topsep=0.07 cm,
        parsep=0.07 cm,
        partopsep=0pt,
        itemsep=0pt,
        leftmargin=0.4 cm + 10pt
    ]
}{
    \end{itemize}
} % new environment for highlights

\newenvironment{highlightsforbulletentries}{
    \begin{itemize}[
        topsep=0.10 cm,
        parsep=0.10 cm,
        partopsep=0pt,
        itemsep=0pt,
        leftmargin=10pt
    ]
}{
    \end{itemize}
} % new environment for highlights for bullet entries


\newenvironment{onecolentry}{
    \begin{adjustwidth}{
        0.2 cm + 0.00001 cm
    }{
        0.2 cm + 0.00001 cm
    }
}{
    \end{adjustwidth}
} % new environment for one column entries

\newenvironment{twocolentry}[2][]{
    \onecolentry
    \def\secondColumn{#2}
    \setcolumnwidth{\fill, 4.5 cm}
    \begin{paracol}{2}
}{
    \switchcolumn \raggedleft \secondColumn
    \end{paracol}
    \endonecolentry
} % new environment for two column entries

\newenvironment{header}{
    \setlength{\topsep}{0pt}\par\kern\topsep\centering\linespread{1.5}
}{
    \par\kern\topsep
} % new environment for the header

\newcommand{\placelastupdatedtext}{% \placetextbox{<horizontal pos>}{<vertical pos>}{<stuff>}
  \AddToShipoutPictureFG*{% Add <stuff> to current page foreground
    \put(
        \LenToUnit{\paperwidth-2 cm-0.2 cm+0.05cm},
        \LenToUnit{\paperheight-1.0 cm})
  }%
}%

% save the original href command in a new command:
\let\hrefWithoutArrow\href

% new command for external links:
\renewcommand{\href}[2]{\hrefWithoutArrow{#1}{\ifthenelse{\equal{#2}{}}{ }{#2 }\raisebox{.15ex}{\footnotesize \faExternalLink*}}}


\begin{document}
    \newcommand{\AND}{\unskip
        \cleaders\copy\ANDbox\hskip\wd\ANDbox
        \ignorespaces
    }
    \newsavebox\ANDbox
    \sbox\ANDbox{}

   
    \begin{header}
        \textbf{\fontsize{24 pt}{24 pt}\selectfont <<FullName>>}

        \vspace{0.1 cm}

        \normalsize
\mbox{
\hrefWithoutArrow{<<LinkedIn>>}{
\color{black}\small\faLinkedinIn\hspace{0.15cm}LinkedIn
}
\kern 0.6cm
\hrefWithoutArrow{mailto:<<Email>>}{
\color{black}\small\faEnvelope[regular]\hspace{0.15cm}<<Email>>
}
\kern 0.6cm
\hrefWithoutArrow{tel:<<Phone>>}{
\color{black}\small\faPhone*\hspace{0.15cm}<<Phone>>
}
\kern 0.6cm
\hrefWithoutArrow{<<GitHub>>}{
\color{black}\small\faGithub\hspace{0.15cm}GitHub
}
\kern 0.6cm
\hrefWithoutArrow{<<LeetCode>>}{
\color{black}\small\faCode\hspace{0.15cm}LeetCode
}
\kern 0.6cm
\hrefWithoutArrow{<<CodeForces>>}{
\color{black}\small\faCode\hspace{0.15cm}CodeForces
}
}
    \end{header}

    \section{Education}
% --- EXAMPLE ENTRY ---
%        \begin{twocolentry}{
%        \textit{2023 – 2027}}
%            \textbf{Netaji Subhas University of Technology}
%            \textit{B.Tech. Electronics \& Communications Engineering (CGPA: 8.38)}
%        \end{twocolentry}
%        \vspace{0.05 cm}
% --- END EXAMPLE ---
<<EducationEntries>>

    \section{Experience}
% --- EXAMPLE ENTRY ---
%        \begin{twocolentry}{
%        \textit{May 2026 – July 2026 }}
%            \textbf{Google}
%            \textit{Intern | }
%            \vspace{0.2 cm}
%            \href{https://github.com/SuryanshKS/protocol-mapper}{\underline{Demo Code}}
%        \end{twocolentry}
%        \begin{onecolentry}
%            \begin{highlights}
%                \item Implemented recursive AST traversal algorithms.
%            \end{highlights}
%        \end{onecolentry}
% --- END EXAMPLE ---
<<ExperienceEntries>>

    
    \section{Projects}
    \vspace{0.05 cm}
% --- EXAMPLE ENTRY ---
%        \textbf{Wandrly - Travel Smarter, Together}\hfill
%        \vspace{0.2 cm}
%\href{https://github.com/SuryanshKS/Wandrly}{\underline{Source Code}}
%|
%\href{https://wandrly-alpha.vercel.app/}{\underline{Live Link}}
%
%        \begin{onecolentry}
%            \begin{highlights}
%                \item Architected a scalable RESTful API backend.
%            \end{highlights}
%        \end{onecolentry}
% --- END EXAMPLE ---
<<ProjectsEntries>>

       
\section{Technical Skills}
\vspace{0.2 cm}
    \begin{highlights}
% --- EXAMPLE ENTRY ---
%        \item Languages: C/C++, JavaScript, Python, SQL
% --- END EXAMPLE ---
<<SkillsEntries>>
    \end{highlights}

 \section{Academic Achievements}        
\vspace{0.1 cm}
\begin{onecolentry}
\begin{highlights}
% --- EXAMPLE ENTRY ---
%    \item Represented my college at the Smart India Hackathon (SIH) 2025 Finals in Nagpur
% --- END EXAMPLE ---
<<AchievementsEntries>>
\end{highlights}
\end{onecolentry}

\section{Extra-Curricular}
    \vspace{0.10 cm}
    \begin{highlights}
% --- EXAMPLE ENTRY ---
%         \item Content Writer | The Debugging Society, NSUT
% --- END EXAMPLE ---
<<ResponsibilityEntries>>
    \end{highlights}
    
\end{document}
